\input{technique-note.tex}
\title{Integration Bee --- Unit 16:\\ Trig Identities \& Intermediate Trig-Sub}
\author{UK Uni Integration Bee prep\\ Yang Topics Guide 2nd Ed §§2.6 Trig-Identities Manipulation, 2.7 Intermediate Trig-Substitution}
\date{}
\begin{document}
\maketitle

\section*{Status}
Yang Chapter 2 concept unit (easy drills only; no past papers).
Tags: \texttt{yang:2.6}, \texttt{yang:2.7}, \texttt{trig-identities}, \texttt{intermediate-trig-sub}.

\section{Official anchors}
\begin{itemize}
  \item Yang Topics Guide 2nd Ed §§2.6 Trig-Identities Manipulation, 2.7 Intermediate Trig-Substitution
        (\texttt{sources/user/Integration\_Bee\_Topics\_Guide\_Yang\_2nd\_Ed.pdf}).
  \item Syllabus §1 techniques support
        (\texttt{sources/UK\_University\_Integration\_Bee\_Syllabus.pdf}).
\end{itemize}

\section{Learning goals}
\begin{enumerate}
  \item Apply power-reduction and product-to-sum identities before integrating.
  \item Use $x=a\sin\theta$ / $x=a\tan\theta$ on intermediate radicals.
  \item Recover circle-area forms like $\int_{0}^{2}\sqrt{4-x^{2}}\,dx=\pi$.
\end{enumerate}

\section{Pedagogy}
Worked $\to$ faded $\to$ interleaved easy drills. Chat Flow 1.

\section{Worked example --- \texttt{U16-D01} (\texttt{yang:2.6})}
\[
\int_{0}^{\pi/2}\sin^{2}x\cos^{2}x\,dx.
\]
$\sin^{2}\cos^{2}=(\sin 2x)^{2}/4$; integrate to $\pi/16$.

\section{Faded example}
Try U16-D05: $\int_{0}^{2}\sqrt{4-x^{2}}\,dx$ (quarter disk of radius 2).

\section{Drill pointers}
See \texttt{drills.tex} / \texttt{drills.md}. SymPy feed: \texttt{verify/drills\_unit16.json}.

\end{document}
